\section{A finite correspondence test for Coulomb relations}
\label{sec:correspondences}

We work with the finite-dimensional algebraic group $\mathsf K=G\times F$.
All homology in this section is equivariant Borel--Moore homology, with
complex coefficients.  We use constructible sheaves only on finite-dimensional
varieties.  The purpose of the section is to express a coefficient relation
in the Coulomb algebra as a relation between such sheaves.

\subsection{Bounded lattice families}

A point $b=gG(\OO)$ of $\Gr_G$ determines the lattice $gN(\OO)$ in
$N(\KK)$.  For a BFN triple $[g,s_0]$, we use $n=gs_0$ as section
coordinate; the fiber over $b$ is $N(\OO)\cap gN(\OO)$.
Let $B$ be a projective $\mathsf K$-variety equipped with a
$\mathsf K$-equivariant map to $\Gr_G$ with bounded image.  The map need not
be an embedding.  For any finite collection of these families, choose $l$
large enough that all their lattices contain $t^lN(\OO)$ and are contained
in $t^{-l}N(\OO)$.  Set
\begin{equation}
 \label{eq:finite-lattice-space}
 \begin{aligned}
 L&=t^{-l}N(\OO)/t^lN(\OO),&
 E_b&=gN(\OO)/t^lN(\OO),\\
 S_B&=\{(b,n)\in B\times L:n\in E_b\},&
 f_B&:S_B\longrightarrow L.
 \end{aligned}
\end{equation}
Here and below the same notation is used for a lattice and for its pullback
to a parameter space.  The representation map
$\Gr_G\to\Gr_{\GL(N)}$ identifies $S_B$ with the pullback of a
tautological vector bundle on a finite Grassmannian.  We work separately on
components if its rank varies.  If $\chi_N=\det(N|_G)$, its rank on a
component of degree $\gamma\in\pi_1(G)$ is
\begin{equation}
 \label{eq:lattice-rank}
 r_B=l\dim N-\langle\chi_N,\gamma\rangle.
\end{equation}
The incidence $S_B\subset B\times L$ is closed, so $f_B$ is proper.
We write $S_{\mathrm{pt}}=N(\OO)/t^lN(\OO)$ for the family at the
identity lattice and $i=f_{\mathrm{pt}}$.  Put $r=l\dim N$ and
$\mathcal F_B=Rf_{B,*}\C_{S_B}$; in particular,
$\mathcal F_{\mathrm{pt}}=i_*\C_{S_{\mathrm{pt}}}$.

Let $Z$ be a finite union of affine Schubert varieties in the neutral
component, containing the identity lattice.  Restriction to that lattice
gives a degree-zero morphism
\begin{equation}
 \label{eq:restriction-at-identity}
 \rho_Z:\mathcal F_Z\longrightarrow\mathcal F_{\mathrm{pt}}.
\end{equation}
Both section bundles here have rank $r$.  The base $Z$ is allowed to be
singular.

The notation $\Hom^d_{\mathsf K,L}(\mathcal F,\mathcal G)$ denotes
$\Hom(\mathcal F,\mathcal G[d])$ for $\mathsf K$-equivariant constructible
complexes on $L$, computed on the common Borel approximations used below.

\begin{lemma}[Finite correspondence factorization]
 \label{lem:correspondence-factorization}
Let $u\in R$ be homogeneous of cohomological degree $d$, and suppose that
in $A_F$ it has an expression
\begin{equation}
 \label{eq:neutral-word-relation}
 u=\sum_w M_{\lambda_{w,1},p_{w,1}}\cdots
                 M_{\lambda_{w,k_w},p_{w,k_w}},
 \qquad k_w\geq1,
 \qquad \sum_j[\lambda_{w,j}]=0\in\pi_1(G).
\end{equation}
Every factor is a nontrivial dressed positive or negative minuscule
generator of Proposition~\ref{prop:minuscule}, and every word is homogeneous
of degree $d$.  Coefficients in $R$ may be included in the dressings.
Choose an ordering of the factors of each word.  There are $Z$ as above
bounding their final lattices and a common truncation
\eqref{eq:finite-lattice-space} such that
\begin{equation}
 \label{eq:finite-relation-factorization}
 u\rho_Z=\sum_w\Psi_w
 \quad\text{in}\quad
 \Hom^d_{\mathsf K,L}
       (\mathcal F_Z,\mathcal F_{\mathrm{pt}}).
\end{equation}
Each $\Psi_w$ factors through a cohomological shift of
$\mathcal F_{\Gr_G^{\lambda_{w,1}}}$.
Only the endpoints are required to lie in $Z$ or in the neutral component;
all intermediate lattices are included in the common truncation.
\end{lemma}

\begin{proof}
We first specify the correspondences and their degrees, then prove the
identity by a faithful localization calculation.

\paragraph{Correspondences and shifts.}
For proper maps $f:X\to L$ and $g:Y\to L$, write $\omega_Y$ for the
complex-linear dualizing complex of $Y$.  Proper base change and adjunction
give
\[
 R\Hom_L(Rf_*\C_X,Rg_*\omega_Y)
 \simeq R\Gamma(X\times_LY,\omega_{X\times_LY}).
\]
For example, if $q:X\times_LY\to Y$ is the projection, the left side is
$R\Hom_Y(q_*\C_{X\times_LY},\omega_Y)$, which gives the displayed
identity by $q_*\dashv q^!$.  Thus, when $Y$ is smooth of complex dimension
$e$, we have
\begin{equation}
 \label{eq:finite-correspondence-Hom}
 \Hom^a_{\mathsf K,L}(Rf_*\C_X,Rg_*\C_Y)
 \simeq H^{\mathsf K}_{2e-a}(X\times_LY).
\end{equation}
The equivariant version follows on common Borel approximations; their
additional dimensions cancel in the shift.  These finite-type models have
triangulable analytic topology and finite cohomological dimension.
We use the proper-support base-change and adjunction formalism of
\cite[Theorems~5.10 and 7.7]{SchnurerSoergel16}, with complex coefficients,
and the duality and Borel--Moore identifications of
\cite[Example~2.7(3) and Corollary~2.11]{MaximSchurmann21}.
In particular, the source
$X$ need not be smooth.  Composition across a smooth middle space is
refined intersection along its diagonal followed by proper pushforward.
These operations are the adjunction and base-change operations in the
displayed identification.

Fix one word, temporarily suppressing its index $w$.  Let $B_j$ be the
variety of chains of lattices with successive relative positions
$\lambda_1,\ldots,\lambda_j$, starting at the identity.  Set
$B_0=\mathrm{pt}$.  Each forgetful map $B_j\to B_{j-1}$ is a bundle with
fiber the closed minuscule orbit $\Gr_G^{\lambda_j}$, so the $B_j$ are
smooth and projective.  Write
\[
 a_j=\dim\Gr_G^{\lambda_j},\qquad
 r_j=\rk E_{b_j},\qquad
 I_j=\{(b_0,\ldots,b_j,n):n\in E_{b_{j-1}}\cap E_{b_j}\}.
\]
The intersection has constant rank by fixed relative position.  It defines
a vector bundle on $B_j$ and a correspondence
\[
 S_{B_j}\ \overset{v_j}{\hookleftarrow}\ I_j
       \xrightarrow{u_j} S_{B_{j-1}}
\]
over $L$, where $v_j$ is the closed subbundle inclusion and $u_j$ forgets
the last lattice.  Write
$c_j=\rk(E_{b_j}/I_j)$, using the same symbol $I_j$ for the intersection
bundle.  The dressing $p_j$ gives an equivariant characteristic class on
$B_j$: use the two tautological blocks of the minuscule modification in
the preceding lattice frame, together with the flavor classes.  This is
the twisted-orbit version of the stabilizer-invariant dressing in
\eqref{eq:minuscule}.

If $\delta_j=\deg p_j$, the class $p_j\cap[I_j]$ defines a morphism
$\mathcal F_{B_j}\to\mathcal F_{B_{j-1}}[q_j]$, where
\begin{equation}
 \label{eq:chain-shift}
 q_j=2(r_{j-1}-r_j+c_j-a_j)+\delta_j.
\end{equation}
Indeed, $\dim I_j=\dim B_{j-1}+a_j+r_j-c_j$, and
\eqref{eq:finite-correspondence-Hom} uses the target dimension
$\dim S_{B_{j-1}}=\dim B_{j-1}+r_{j-1}$.
The endpoint map $B_k\to Z$ induces a proper map $S_{B_k}\to S_Z$.
Its pullback morphism on constant sheaves, followed by pushforward to
$L$, gives $\mathcal F_Z\to\mathcal F_{B_k}$ with no shift.
Compose it with the $k$ incidence morphisms to obtain $\Psi_w$.
This already gives the stated first-step factorization.

At a modification of type $\lambda_j$,
$c_j=\sum_\beta(-\langle\beta,\lambda_j\rangle)_+$, with multiplicities
and with the pairing taken on the gauge part.  The rational coefficient
in Equation~\eqref{eq:minuscule} has degree $2c_j-2a_j+\delta_j$,
counting linear forms with degree two.  For the neutral word, the sum of
these degrees is its BFN degree $d$, by
Proposition~\ref{prop:minuscule}.  Moreover,
Equation~\eqref{eq:lattice-rank} gives $r_k=r_0=r$.  Consequently
\[
 \sum_jq_j
 =2(r_0-r_k)+\sum_j(2c_j-2a_j+\delta_j)=d.
\]
Thus $\Psi_w$ and $u\rho_Z$ belong to the same graded Hom space.
No equality of intermediate ranks or identification of nonneutral degree
conventions is used.

\paragraph{Faithfulness of the localization test.}
Put $C_Z=S_{\mathrm{pt}}\times_LS_Z$ and let $\iota:C_Z\hookrightarrow
S_Z$ be the inclusion.  The identity
\eqref{eq:finite-correspondence-Hom}, with smooth target
$S_{\mathrm{pt}}$, identifies the Hom space in
\eqref{eq:finite-relation-factorization} with
$H^{\mathsf K}_{2r-d}(C_Z)$.  Under it, $\rho_Z$ is the fundamental class
of $S_{\mathrm{pt}}$ over the identity lattice.

Let $P=H^*_{T\times F}(\mathrm{pt})$.  The stratification of $C_Z$ over
the finitely many $G(\OO)$-orbits in $Z$ has vector-bundle strata.  Each
orbit is an affine-space bundle over a partial flag variety and contracts
to it under loop scaling $t\mapsto at$ as $a\to0$.  This contraction
commutes with $T$, and $F$ acts trivially on the orbit.  Its ordinary
equivariant cohomology is therefore even and free over $P$; smooth
equivariant Poincar\'e duality gives the same assertion for Borel--Moore
homology, with a dimension shift.  The section bundle adds an even Thom
shift.
The closed-stratum exact sequences split into short exact sequences by
parity, and their free quotients make each extension free over $P$.
It follows that $H_*^{T\times F}(C_Z)$ is free, in particular
torsion-free.  This is the finite version of the orbit-filtration argument
in \cite[Lemmas~2.5--2.6]{BFN18}.

The $T$-fixed lattices in $\Gr_G$ are precisely the diagonal lattices
$t^\eta$: a fixed lattice in each $V_i(\KK)$ splits along its distinct
diagonal characters, and a lattice in a one-dimensional Laurent series
space is $t^m\OO$.  At $t^\eta$, every $(T\times F)$-weight-zero
summand of $N$ has gauge pairing zero with $\eta$.  Its lattice is
unchanged by intersection with $N(\OO)$.  Thus $\iota$ identifies the
fixed loci of $C_Z$ and $S_Z$.  Equivariant localization makes
$\iota_*$ an isomorphism after tensoring with $\operatorname{Frac}P$;
torsion-freeness makes this localized pushforward a faithful test on
$H_*^{T\times F}(C_Z)$.  Passing to Weyl invariants recovers
$\mathsf K$-equivariance.  This is also the fixed-fiber calculation of
\cite[Lemma~5.13]{BFN18}.

After localization, the classes of the entire section fibers
$E_{t^\eta}$ form a basis for the homology of $S_Z$, indexed by its
fixed lattices.  These classes, rather than isolated points in a fiber,
allow weight-zero summands of $N$ without inverting a zero Euler class.
We denote this basis formally by $e^\eta$.

\paragraph{The chain calculation.}
For each partial chain put $C_j=S_{\mathrm{pt}}\times_LS_{B_j}$, and let
$h_j:C_j\hookrightarrow S_{B_j}$ forget the initial section constraint.
Let $z_j$ be the Borel--Moore class of the composite
$\mathcal F_{B_j}\to\mathcal F_{\mathrm{pt}}[\sum_{h=1}^j q_h]$
formed by the first $j$ incidence morphisms.  Since $S_{B_j}$ is smooth, write
$(h_j)_*z_j=\alpha_j\cap[S_{B_j}]$.  We claim that
\begin{equation}
 \label{eq:chain-euler-product}
 \alpha_k=\prod_{j=1}^k p_j\,\eu(E_{b_j}/I_j),
\end{equation}
with all factors pulled to $B_k$.

The initial class has $\alpha_0=1$.  At the $j$th step, composition is
refined intersection along the diagonal of the smooth middle space
$S_{B_{j-1}}$.  This intersection commutes with the proper pushforward
that forgets the initial section constraint.  We may therefore first
replace $z_{j-1}$ by $(h_{j-1})_*z_{j-1}
=\alpha_{j-1}\cap[S_{B_{j-1}}]$.  The remaining intersection pulls this
cohomology class back to $I_j$, giving
\[
 (h_j)_*z_j
 =(v_j)_*\bigl(u_j^*\alpha_{j-1}\,p_j\cap[I_j]\bigr).
\]
Here the regular diagonal belongs to the smooth middle space; no
smoothness of $u_j$ is needed.  The classes $\alpha_{j-1}$ are pulled
from the chain bases, and the subbundle inclusion $v_j$ multiplies by
$\eu(E_{b_j}/I_j)$.  This proves \eqref{eq:chain-euler-product} by
induction.

Under the same Hom--Borel--Moore identification, precomposition with the
unit for the proper endpoint map $S_{B_k}\to S_Z$ is proper pushforward
along $C_k\to C_Z$.  Forgetting the initial section constraint commutes
with this pushforward.  Thus the endpoint pushforward of
$\alpha_k\cap[S_{B_k}]$ is exactly the image of $\Psi_w$ in the
localization test.

Fixed chains are sequences of Weyl translates $\lambda'_j$ of the chosen
minuscule coweights; their successive lattices are
$t^{\eta_j}$ with $\eta_j=\sum_{h\leq j}\lambda'_h$.
At such a chain the $j$th modification is translated by
$t^{\eta_{j-1}}$.  This translation commutes with the diagonal torus.
There is no loop rotation, so it changes neither the relevant weights nor
the dressing values.  Hence its quotient Euler class is
\[
 \prod_\beta\beta^{(-\langle\beta,\lambda'_j\rangle)_+},
\]
and its relative tangent Euler class is
$\eu_T(T_{t^{\lambda'_j}}\Gr_G^{\lambda_j})$.
Localization on the smooth chain base therefore sends
\eqref{eq:chain-euler-product}, after endpoint pushforward, to
\[
 \sum_{(\lambda'_1,\ldots,\lambda'_k)}
 e^{\lambda'_1+\cdots+\lambda'_k}
 \prod_{j=1}^k
 \left(
 p_{j,\lambda'_j}
 \frac{\prod_\beta\beta^{(-\langle\beta,\lambda'_j\rangle)_+}}
 {\eu_T(T_{t^{\lambda'_j}}\Gr_G^{\lambda_j})}
 \right).
\]
This is exactly the product of the expressions
\eqref{eq:minuscule} for the word.  All dressings are transported with
their coweights.  The fiber ranks at the endpoint equal $r$, so no
additional fiber factor enters endpoint pushforward.

The localized pushforward of $u\rho_Z$ is $u e^0$.
Equation~\eqref{eq:neutral-word-relation} and the abelianization of
Proposition~\ref{prop:minuscule} identify it with the sum of the displayed
word expressions.  The faithful test just proved implies the equality
in $H^{\mathsf K}_{2r-d}(C_Z)$ and hence
\eqref{eq:finite-relation-factorization}.
\end{proof}

\begin{remark}
Every $u\in\ker(R\to B_F)$ admits an expression of the form
\eqref{eq:neutral-word-relation}, degree by degree.  Indeed,
Proposition~\ref{prop:minuscule} implies that the ideal of all nonzero
topological degrees is generated by the nontrivial positive and negative
minuscule generators.  Expand a relation in that ideal and retain its
neutral component.  Commutativity of $A_F$ permits any desired ordering
of the factors before applying Lemma~\ref{lem:correspondence-factorization}.
\end{remark}

The relation has now been transferred to a finite-dimensional category,
with the identity-lattice map on one side and first-step factorizations
on the other.  The stable-pair test in Section~\ref{sec:stable-pairs}
will annihilate the selected first-step objects while keeping a split
restriction from the identity-lattice test onto cohomology.
